\editionPage[\pdfbookmark[0]{Deux bilans annuels}{bm-bilans-guide-parent}]{bilans-guide}{Deux bilans pour l’année}{Sujets originaux · Relier les chapitres}{Durée conseillée : 1 h 30 · Barème d’entraînement sur 20}
\begin{blocEdition}{bilans-guide}

Ces évaluations sont des créations du livre pour la Sixième. Elles n’ont
pas de statut d’examen officiel. Le temps et le barème sont des repères pour
une séance autonome ; ils s’adaptent librement à la classe ou à la maison.
\editionEtape{Automatismes : environ 20 minutes}
Douze questions courtes, sans calculatrice ni brouillon long : 6 points,
soit 0,5 point par réponse. Écrire les étapes utiles, même brèves.
\editionEtape{Problèmes : environ 1 h 10}
Quatre problèmes indépendants : 14 points. Chaque problème vaut
3,5 points : trois questions à 1 point et 0,5 point pour la clarté de la
rédaction, les unités et les contrôles.
\editionEtape{Travailler après l’évaluation}
Comparer les démarches, relever les erreurs par notion, puis refaire
les questions utiles sans correction. Une démarche peut être juste
même quand le résultat final est faux : repérer l’étape à reprendre.
\editionEtape{Préparer une autre séance}
Revenir au cours et à deux exercices du chapitre concerné. Le second
bilan peut être fait plus tard pour vérifier que la méthode est acquise.

\end{blocEdition}
\editionPage{bilan-A-auto}{Bilan A}{Nombres, calculs et proportions}{Partie 1 · 20 min · Sans calculatrice · 6 points}
\begin{blocEdition}{bilan-A-auto}
\begin{questions}
\item Écrire en chiffres : « deux-milliards-quarante-millions ».
\item Écrire $3+\frac5{10}+\frac7{1\,000}$ avec une virgule.
\item Ranger dans l’ordre croissant : $0{,}7$ ; $0{,}65$ ; $0{,}699$.
\item Arrondir $14{,}386$ au centième.
\item Calculer $2{,}5\times0{,}4$.
\item Calculer $56\times0{,}01$.
\item Quotient et reste de la division euclidienne de 200 par 15.
\item Compléter : $7\times\square=5$.
\item Calculer $\frac34-\frac14$, puis $\frac12+\frac14$.
\item Calculer les $\frac23$ de 18.
\item Calculer 25~\% de 36.
\item 3 stylos coûtent $4{,}50$ €. Combien coûtent 6 stylos ?\end{questions}
\end{blocEdition}
\editionPage{bilan-A-ex-1}{Bilan A · Problème 1}{La médiathèque}{Partie 2 · 1 h 10 pour les quatre problèmes · 3,5 points}
\begin{blocEdition}{bilan-A-ex-1}
\begin{questions}
\item Une médiathèque possède $12\,480$ livres. Elle en achète 350 par an. Combien en aura-t-elle dans 4 ans, sans en retirer ?
\item Elle range ses $12\,480$ livres sur des étagères de 48 livres. Combien d’étagères pleines ? Reste-t-il des livres ?
\item Un livre neuf coûte en moyenne $18{,}50$ €. Combien coûtent les 350 livres achetés en un an ?
\end{questions}\par\vspace{5mm}{\otLight\fontsize{8}{10}\selectfont\color{otMuted}Recherche et rédaction}\par\par\vspace{7mm}{\color{otRule!65}\hrule height .25pt}\par\vspace{7mm}{\color{otRule!65}\hrule height .25pt}\par\vspace{7mm}{\color{otRule!65}\hrule height .25pt}\par\vspace{7mm}{\color{otRule!65}\hrule height .25pt}\par\vspace{7mm}{\color{otRule!65}\hrule height .25pt}\par\vspace{7mm}{\color{otRule!65}\hrule height .25pt}
\end{blocEdition}
\editionPage{bilan-A-ex-2}{Bilan A · Problème 2}{Le marché}{Partie 2 · 1 h 10 pour les quatre problèmes · 3,5 points}
\begin{blocEdition}{bilan-A-ex-2}
\begin{questions}
\item Des clémentines sont vendues $2{,}80$ € le kilogramme. Combien coûtent $2{,}5$ kg ?
\item Les cerises coûtent 6 € pour 750 g. Quel est leur prix au kilogramme ?
\item Inès paie avec un billet de 20 € ses $2{,}5$ kg de clémentines et 500 g de cerises. Combien lui rend-on ?
\end{questions}\par\vspace{5mm}{\otLight\fontsize{8}{10}\selectfont\color{otMuted}Recherche et rédaction}\par\par\vspace{7mm}{\color{otRule!65}\hrule height .25pt}\par\vspace{7mm}{\color{otRule!65}\hrule height .25pt}\par\vspace{7mm}{\color{otRule!65}\hrule height .25pt}\par\vspace{7mm}{\color{otRule!65}\hrule height .25pt}\par\vspace{7mm}{\color{otRule!65}\hrule height .25pt}\par\vspace{7mm}{\color{otRule!65}\hrule height .25pt}
\end{blocEdition}
\editionPage{bilan-A-ex-3}{Bilan A · Problème 3}{Le gâteau et les parts}{Partie 2 · 1 h 10 pour les quatre problèmes · 3,5 points}
\begin{blocEdition}{bilan-A-ex-3}
\begin{questions}
\item Un gâteau est partagé : Léa en prend $\frac14$, Noé $\frac13$. Quelle fraction du gâteau ont-ils prise ensemble ?
\item Quelle fraction reste-t-il ? L’écrire avec le dénominateur 12.
\item Le gâteau pesait 840 g. Quelle masse reste-t-il ? Quel pourcentage du gâteau Léa a-t-elle mangé ?
\end{questions}\par\vspace{5mm}{\otLight\fontsize{8}{10}\selectfont\color{otMuted}Recherche et rédaction}\par\par\vspace{7mm}{\color{otRule!65}\hrule height .25pt}\par\vspace{7mm}{\color{otRule!65}\hrule height .25pt}\par\vspace{7mm}{\color{otRule!65}\hrule height .25pt}\par\vspace{7mm}{\color{otRule!65}\hrule height .25pt}\par\vspace{7mm}{\color{otRule!65}\hrule height .25pt}\par\vspace{7mm}{\color{otRule!65}\hrule height .25pt}
\end{blocEdition}
\editionPage{bilan-A-ex-4}{Bilan A · Problème 4}{Les cartes de collection}{Partie 2 · 1 h 10 pour les quatre problèmes · 3,5 points}
\begin{blocEdition}{bilan-A-ex-4}
\begin{questions}
\item Tom a 3 fois plus de cartes que Sarah ; ensemble, ils en ont 96. Représenter la situation par un schéma en barres, puis trouver le nombre de cartes de chacun.
\item Tom donne 12 cartes à Sarah. Combien chacun en a-t-il alors ?
\item Combien Tom aurait-il dû en donner pour qu’ils en aient autant ?
\end{questions}\par\vspace{5mm}{\otLight\fontsize{8}{10}\selectfont\color{otMuted}Recherche et rédaction}\par\par\vspace{7mm}{\color{otRule!65}\hrule height .25pt}\par\vspace{7mm}{\color{otRule!65}\hrule height .25pt}\par\vspace{7mm}{\color{otRule!65}\hrule height .25pt}\par\vspace{7mm}{\color{otRule!65}\hrule height .25pt}\par\vspace{7mm}{\color{otRule!65}\hrule height .25pt}\par\vspace{7mm}{\color{otRule!65}\hrule height .25pt}
\end{blocEdition}
\editionPage{bilan-B-auto}{Bilan B}{Grandeurs, géométrie et données}{Partie 1 · 20 min · Sans calculatrice · 6 points}
\begin{blocEdition}{bilan-B-auto}
\begin{questions}
\item Convertir $3{,}5$ km en mètres.
\item Convertir 450 cm$^2$ en dm$^2$.
\item Aire d’un rectangle de 8 m sur $2{,}5$ m.
\item Périmètre d’un disque de diamètre 20 cm, avec $\pi\approx3{,}14$.
\item Écrire $2{,}25$ h en heures et minutes.
\item Un triangle a deux angles de $50^\circ$ et $60^\circ$. Mesure du troisième ?
\item Un triangle équilatéral : mesure de chaque angle ?
\item Deux angles supplémentaires : l’un mesure $35^\circ$. L’autre ?
\item Le point M est sur la médiatrice de [AB] et $MA=6$ cm. Que vaut MB ?
\item Combien d’axes de symétrie a un carré ?
\item Volume d’un pavé de 3 cubes sur 2 sur 2, en cubes de 1 cm$^3$.
\item On lance un dé équilibré. Probabilité d’obtenir un multiple de 3 ?\end{questions}
\end{blocEdition}
\editionPage{bilan-B-ex-1}{Bilan B · Problème 1}{La cour du collège}{Partie 2 · 1 h 10 pour les quatre problèmes · 3,5 points}
\begin{blocEdition}{bilan-B-ex-1}
\begin{questions}
\item La cour est un rectangle de 45 m sur 30 m. Calculer son périmètre et son aire.
\item On y trace une piste circulaire de 20 m de diamètre. Quelle est la longueur d’un tour, au dixième de mètre ?
\item Un élève fait 6 tours de piste en 3 min 30 s. Quelle distance parcourt-il ? Combien de temps met-il en moyenne pour un tour ?
\end{questions}\par\vspace{5mm}{\otLight\fontsize{8}{10}\selectfont\color{otMuted}Recherche et rédaction}\par\par\vspace{7mm}{\color{otRule!65}\hrule height .25pt}\par\vspace{7mm}{\color{otRule!65}\hrule height .25pt}\par\vspace{7mm}{\color{otRule!65}\hrule height .25pt}\par\vspace{7mm}{\color{otRule!65}\hrule height .25pt}\par\vspace{7mm}{\color{otRule!65}\hrule height .25pt}\par\vspace{7mm}{\color{otRule!65}\hrule height .25pt}
\end{blocEdition}
\editionPage{bilan-B-ex-2}{Bilan B · Problème 2}{Le triangle et son cercle}{Partie 2 · 1 h 10 pour les quatre problèmes · 3,5 points}
\begin{blocEdition}{bilan-B-ex-2}
\begin{questions}
\item Construire le triangle ABC tel que $AB=7$ cm, $\widehat{A}=40^\circ$ et $\widehat{B}=60^\circ$. Calculer $\widehat{C}$.
\item Construire le cercle circonscrit au triangle. Expliquer la construction.
\item Construire le symétrique $A'B'C'$ du triangle par rapport à la droite (BC). Que vaut $A'B'$ ? L’angle de sommet $B'$ ?
\end{questions}\label{fig-bilan-B-ex-2}\figureGeo{bilan-triangle}
\par\vspace{5mm}{\otLight\fontsize{8}{10}\selectfont\color{otMuted}Recherche et rédaction}\par\par\vspace{7mm}{\color{otRule!65}\hrule height .25pt}\par\vspace{7mm}{\color{otRule!65}\hrule height .25pt}\par\vspace{7mm}{\color{otRule!65}\hrule height .25pt}\par\vspace{7mm}{\color{otRule!65}\hrule height .25pt}\par\vspace{7mm}{\color{otRule!65}\hrule height .25pt}\par\vspace{7mm}{\color{otRule!65}\hrule height .25pt}
\end{blocEdition}
\editionPage{bilan-B-ex-3}{Bilan B · Problème 3}{La sortie au zoo}{Partie 2 · 1 h 10 pour les quatre problèmes · 3,5 points}
\begin{blocEdition}{bilan-B-ex-3}
\begin{questions}
\item Le car part du collège à 8 h 40 et arrive au zoo à 10 h 15. Quelle est la durée du trajet ?
\item La visite dure 3,5 h à partir de l’arrivée. À quelle heure se termine-t-elle ?
\item Sur le plan du zoo, 1 cm représente 50 m. L’allée des fauves mesure 9 cm sur le plan. Quelle est sa longueur réelle ? Combien de temps faut-il pour la parcourir à 60 m par minute ?
\end{questions}\par\vspace{5mm}{\otLight\fontsize{8}{10}\selectfont\color{otMuted}Recherche et rédaction}\par\par\vspace{7mm}{\color{otRule!65}\hrule height .25pt}\par\vspace{7mm}{\color{otRule!65}\hrule height .25pt}\par\vspace{7mm}{\color{otRule!65}\hrule height .25pt}\par\vspace{7mm}{\color{otRule!65}\hrule height .25pt}\par\vspace{7mm}{\color{otRule!65}\hrule height .25pt}\par\vspace{7mm}{\color{otRule!65}\hrule height .25pt}
\end{blocEdition}
\editionPage{bilan-B-ex-4}{Bilan B · Problème 4}{Les jeux de la kermesse}{Partie 2 · 1 h 10 pour les quatre problèmes · 3,5 points}
\begin{blocEdition}{bilan-B-ex-4}
\begin{questions}
\item Une urne contient 4 boules rouges, 6 boules vertes et 10 boules jaunes. Quelle est la probabilité de tirer une boule verte ? L’écrire aussi en pourcentage.
\item Le tableau donne le nombre de lots gagnés à chaque stand. Quel stand a distribué le plus de lots ? Combien de lots en tout ?\begin{center}\renewcommand{\arraystretch}{1.35}\begin{tabular}{c|cccc}\toprule
Stand & pêche & chamboule-tout & loterie & quiz\\\midrule
Lots & 35 & 28 & 42 & 15\\
\bottomrule\end{tabular}\end{center}
\item On retire de l’urne toutes les boules jaunes. Quelle est la nouvelle probabilité de tirer une boule verte ?
\end{questions}\par\vspace{5mm}{\otLight\fontsize{8}{10}\selectfont\color{otMuted}Recherche et rédaction}\par\par\vspace{7mm}{\color{otRule!65}\hrule height .25pt}\par\vspace{7mm}{\color{otRule!65}\hrule height .25pt}\par\vspace{7mm}{\color{otRule!65}\hrule height .25pt}\par\vspace{7mm}{\color{otRule!65}\hrule height .25pt}\par\vspace{7mm}{\color{otRule!65}\hrule height .25pt}\par\vspace{7mm}{\color{otRule!65}\hrule height .25pt}
\end{blocEdition}
